% LAB BOREDOM — module L1, The Three Literatures. FCDP Stage D2, batch 1 of 6.
% Pack: lab-boredom-L1-pack-v2 (+ -v2-quotes.json)  |  Plan: L1-D1-MOVEMENT-PLAN-v2-2026-08-06
% Skeleton: L1-D15-SKELETON-v1-2026-08-06  |  Model: Fable 5 @ xhigh  |  No retrieval in drafting context.
% Corrections carried: Q01 BARRED (abstract-only); Q02 cited at printed 484; Q05 attributed to Fahlman.
% Batch 1 revision r1 (author, 2026-08-06): opener rephrased; twenty-four-year sentence deleted;
%   "fifty-five works" -> "the works surveyed here"; E18 footnote deleted (E18 re-placed to P12);
%   King & Salvo studies moved to FIRST PERSON and referred to directly.
% Open in this batch: the self-citing author of E01 is unnamed — Stage-P resolution in flight.

Boredom has been academically pursued across two rather different fields of study. One asks
what boredom discloses about the person who suffers it, treats the experience as a disclosure
of how things stand for someone rather than as a state to be relieved, and reaches for
Heidegger and the phenomenological tradition that follows him; the other asks what boredom
looks like in a signal, and reaches for electrodes, eye trackers, and self-report scales
validated against one another. Across the works surveyed here, the connection between them
amounts to a single author citing their own work across two programs, together with one
critical footnote. That is a citation trail rather than an exchange. The first counts a
reading of experience as evidence; the second counts a number extracted from a sensor.
Neither treats what the other counts as bearing on its own question, and the absence is not an
oversight either could have corrected by wider reading. The two studies we published stand in
neither literature: they belong to a third lineage, clinical immersion in engineering
education, where a headset substitutes for clinical access that no program can secure in
sufficient quantity, and that lineage arrives with its own question, its own warrant, its own
reference list, and no commitment about what boredom is. From the second literature we took
one thing, deliberately and by name: the definition of boredom under which we collected the
data. That definition reached us through Fahlman from Eastwood, it is the definition our
methods section declares, it governs what our instruments were built to detect, and it is also
the definition that could not name what we found. What can now be read out of this data
follows from that borrowing.
